\RequirePackage{fix-cm}
\documentclass[12pt,a4paper]{article}

\usepackage{pdfpages}
\usepackage{graphicx}
\usepackage{tikz}

\usepackage{fontspec}
\usepackage{polyglossia}

\setmainlanguage[numerals=arabic]{armenian}

\setmainfont{GHEA Grapalat}
\newfontfamily\armenianfont{GHEA Grapalat}
\newfontfamily\armenianfonttt{GHEA Grapalat}

\usepackage[a4paper,
            top=2.5cm,
            bottom=2.5cm,
            left=2.5cm,
            right=2.5cm]{geometry}

\usepackage{setspace}
\usepackage{booktabs}
\usepackage{microtype}
\usepackage{titlesec}
\usepackage{amsmath}
\usepackage{amsfonts}
\usepackage{amssymb}

\numberwithin{equation}{section}
\numberwithin{figure}{section}

\usepackage[hidelinks]{hyperref}
\usepackage[nottoc]{tocbibind}

\titleformat{\section}
{\centering\bfseries\fontsize{14}{16}\selectfont}
{\thesection}{1em}{\MakeUppercase}

\let\oldsection\section
\RenewDocumentCommand{\section}{s o m}{%
  \IfBooleanTF{#1}
    {\oldsection*{#3}}
    {\IfNoValueTF{#2}
      {\oldsection[\texorpdfstring{\MakeUppercase{#3}}{#3}]{#3}}
      {\oldsection[\texorpdfstring{\MakeUppercase{#2}}{#2}]{#3}}}%
}

\titleformat{\subsection}
{\centering\bfseries\fontsize{13}{15}\selectfont}
{\thesubsection}{1em}{}

\setlength{\emergencystretch}{3em}
\setlength{\parindent}{0.75cm}
\setlength{\parskip}{0pt}
\setstretch{1.5}
\sloppy

\usepackage{listings}
\usepackage{xcolor}   

\lstdefinelanguage{Julia}
  {morekeywords={struct,end,function,if,else,for,while,using,module,const,let},
   sensitive=true,
   morecomment=[l]{\#},
   morestring=[b]",
  }[keywords,comments,strings]

\lstset{
    backgroundcolor=\color{gray!5},   
    basicstyle=\ttfamily\small,       
    keywordstyle=\color{blue},        
    commentstyle=\color{gray},        
    stringstyle=\color{purple},       
    showstringspaces=false,
    breaklines=true,                  
}

\usepackage{enumitem}
\setlist{nosep, leftmargin=1.25cm}

\begin{document}

\renewcommand{\contentsname}{ԲՈՎԱՆԴԱԿՈՒԹՅՈՒՆ}
\tableofcontents

\newpage

\renewcommand{\thefigure}{\arabic{figure}}
\section*{Նախաբան}
\phantomsection
\addcontentsline{toc}{section}{ՆԱԽԱԲԱՆ}

Դիցուք հարթության վրա տրված է ինչ-որ կետերի հավաքածու։ Այդ դեպքում այդ հավաքածուն ապրոքսիմացնել ինչ-որ կերպով նշանակում է գտնել ինչ-որ $L$ կոր, որը պարտադիր չէ, որ անցնի տրված բոլոր կետերով, բայց հնարավորինս «լավ» կբնութագրի կետերի այդ հավաքածուն։

\begin{figure}[htbp]
    \centering
    \begin{tikzpicture}[scale=1]
        % Կոորդինատային առանցքներ
        \draw[->] (-1.3,0) -- (4.8,0) node[right] {$x$};
        \draw[->] (0,-0.4) -- (0,2.8) node[left] {$y$};

        % Ապրոքսիմացնող L կոր
        \draw[thick] plot[smooth] coordinates {(-0.9,1.0) (0,1.75) (1.2,1.5) (2.6,1.15) (4.2,1.55)};
        \node[below] at (2.2,1.15) {$L$};

        % Տրված կետեր
        \fill (-0.7,1.35) circle (2pt);
        \fill (0.05,1.9) circle (2pt);
        \fill (1.3,1.85) circle (2pt);
        \fill (2.6,0.9) circle (2pt);
        \fill (4.0,1.8) circle (2pt);
    \end{tikzpicture}
    \caption{Ապրոքսիմացիա}
\end{figure}

Եթե կորն անցնում է տրված բոլոր կետերով, ապա այն կոչվում է տրված կետերի բազմության ինտերպոլյացիոն կոր։

\begin{figure}[htbp]
    \centering
    \begin{tikzpicture}[scale=1]
        % Կոորդինատային առանցքներ
        \draw[->] (-2.2,0) -- (4.8,0) node[right] {$x$};
        \draw[->] (0,-0.4) -- (0,2.8) node[right] {$y$};

        % Ինտերպոլյացիոն L կոր
        \draw[thick] plot[smooth] coordinates {(-1.8,0.7) (-1.5,1.0) (-0.1,1.8) (1.1,1.7) (2.3,1.7) (3.6,2.2) (4.1,2.55)};
        \node[below right] at (3.9,2.4) {$L$};

        % Տրված կետեր
        \fill (-1.5,1.0) circle (2pt);
        \fill (-0.1,1.8) circle (2pt);
        \fill (1.1,1.7) circle (2pt);
        \fill (2.3,1.7) circle (2pt);
        \fill (3.6,2.2) circle (2pt);
    \end{tikzpicture}
    \caption{Ինտերպոլյացիա}
\end{figure}

\renewcommand{\thefigure}{\thesection.\arabic{figure}}

Գործնական խնդիրներում (ֆիզիկայում, տնտեսագիտության մեջ, ինժեներիայում և տվյալների վերլուծության ոլորտում) փորձարարական կամ դիտարկումների միջոցով ստացված տվյալները սովորաբար պարունակում են չափման պատահական սխալներ և աղմուկ։ Այդպիսի տվյալների դեպքում ինտերպոլյացիոն կորի կառուցումը հաճախ աննպատակահարմար է, քանի որ կորը ստիպված է լինում անցնել նաև աղմկոտ կետերով՝ առաջացնելով անցանկալի տատանումներ և խեղաթյուրելով ընդհանուր օրինաչափությունը։ Հետևաբար, անհրաժեշտություն է առաջանում կառուցել համեմատաբար պարզ և հարթ մոտարկող ֆունկցիա, որը չի կրկնի պատահական շեղումները, այլ կարտահայտի տվյալների հիմնական միտումը։

Այդպիսի խնդիրների լուծման ամենակիրառելի և դասական եղանակը \textbf{փոքրագույն քառակուսիների մեթոդն է} (առաջարկված Կ. Ֆ. Գաուսի և Ա. Մ. Լեժանդրի կողմից), որի հիմքում ընկած է դիտարկվող կետերից մոտարկող կորի շեղումների քառակուսիների գումարի նվազեցումը։

Սույն աշխատանքի նպատակն է ուսումնասիրել փոքրագույն քառակուսիների մեթոդի մաթեմատիկական հիմունքները գծային և բազմանդամային ապրոքսիմացիայի խնդիրներում։ Աշխատանքում դիտարկվում են գծային և $m$-րդ աստիճանի բազմանդամային ֆունկցիաների կառուցումը, լուծման գոյության ու միակության պայմանները, մատրիցային ձևակերպումը, մոտարկման ճշգրտության գնահատման չափանիշները, ինչպես նաև կոնկրետ թվային օրինակների վերլուծությունը։

\section[գծային ապրոքսիմացիա \texorpdfstring{\newline}{}(փոքրագույն քառակուսիների մեթոդով)]{գծային ապրոքսիմացիա (փոքրագույն քառակուսիների մեթոդով)}

Դիցուք $Oxy$ կոորդինատական հարթությունում տրված են $M_1(x_1, y_1)$, $M_2(x_2, y_2)$, $\dots$, $M_n(x_n, y_n)$ կետերը ($n \geqslant 3$)։

\begin{figure}[htbp]
    \centering
    \begin{tikzpicture}[scale=1]
        % Կոորդինատային առանցքներ
        \draw[->] (-3.4,0) -- (5.4,0) node[right] {$x$};
        \draw[->] (0,-0.6) -- (0,3.5) node[right] {$y$};
        \node[below left] at (0,0) {$O$};

        % y = kx + b ուղիղը (y = 0.35x + 1.0)
        \draw[thick] (-3.2,-0.12) -- (5.0,2.75) node[right] {$y = kx + b$};

        % M_1 կետ
        \fill (-1.6,1.15) circle (2pt) node[above left] {$M_1$};
        \draw (-1.6,1.15) -- (-1.6,0.44);
        \draw[dashed] (-1.6,0.44) -- (-1.6,0) node[below] {$x_1$};

        % M_2 կետ
        \fill (-0.8,2.7) circle (2pt) node[left] {$M_2$};
        \draw (-0.8,2.7) -- (-0.8,0.72);
        \draw[dashed] (-0.8,0.72) -- (-0.8,0) node[below] {$x_2$};

        % M_3 կետ
        \fill (0.8,2.0) circle (2pt) node[above left] {$M_3$};
        \draw (0.8,2.0) -- (0.8,1.28);
        \draw[dashed] (0.8,1.28) -- (0.8,0) node[below] {$x_3$};

        % M_4 կետ (ուղղից ներքև)
        \fill (1.6,0.75) circle (2pt) node[right] {$M_4$};
        \draw (1.6,1.56) -- (1.6,0.75);
        \draw[dashed] (1.6,0.75) -- (1.6,0) node[below] {$x_4$};

        % M_5 կետ
        \fill (2.4,2.6) circle (2pt) node[above left] {$M_5$};
        \draw (2.4,2.6) -- (2.4,1.84);
        \draw[dashed] (2.4,1.84) -- (2.4,0) node[below] {$x_5$};

        % Միջանկյալ կետ (ուղղից ներքև)
        \fill (3.2,1.45) circle (2pt);
        \draw (3.2,2.12) -- (3.2,1.45);

        % M_n կետ
        \fill (4.2,3.15) circle (2pt) node[above right] {$M_n$};
        \draw (4.2,3.15) -- (4.2,2.47);
        \draw[dashed] (4.2,2.47) -- (4.2,0) node[below] {$x_n$};
    \end{tikzpicture}
    \caption{գծային ապրոքսիմացիա փոքրագույն քառակուսիների մեթոդով}
\end{figure}

Պահանջվում է գտնել $y = kx + b$ գծային ֆունկցիան այնպես, որ այն հնարավորինս «լավ» բնութագրի կետերի տրված բազմությունը։

Դիտարկենք հետևյալ մեծությունը՝
\[
S(k, b) = \sum_{i=1}^n \bigl(y_i - (kx_i + b)\bigr)^2 \to \min
\]
Պետք է գտնել $k$-ն և $b$-ն։

Քանի որ $S(k, b) \geqslant 0$ անընդհատ և դիֆերենցելի ֆունկցիա է, և $S(k, b) \underset{\substack{k \to \infty \\ b \to \infty}}{\longrightarrow} +\infty$, ուստի եթե $S(k, b)$ ֆունկցիան ունի ճիշտ մեկ ստացիոնար կետ, ապա այդ կետում այն կընդունի իր փոքրագույն արժեքը։

Ցույց տանք, որ $S(k, b)$ ֆունկցիան ունի ճիշտ մեկ ստացիոնար կետ։

\begin{align*}
\frac{\partial S}{\partial k} &= -2 \sum_{i=1}^n \bigl(y_i - (kx_i + b)\bigr) x_i = -2 \sum_{i=1}^n \bigl(x_i y_i - (kx_i^2 + b x_i)\bigr) \\
\frac{\partial S}{\partial b} &= -2 \sum_{i=1}^n \bigl(y_i - (kx_i + b)\bigr)
\end{align*}

Կազմենք ստացիոնար կետերի գտնելու համակարգը՝
\begin{align*}
\begin{cases}
\dfrac{\partial S}{\partial b} = 0 \\[8pt]
\dfrac{\partial S}{\partial k} = 0
\end{cases}
&\Longleftrightarrow
\begin{cases}
\displaystyle\sum_{i=1}^n \bigl(y_i - (kx_i + b)\bigr) = 0 \\[12pt]
\displaystyle\sum_{i=1}^n \bigl(x_i y_i - (kx_i^2 + b x_i)\bigr) = 0
\end{cases} \\[10pt]
&\Longleftrightarrow
\begin{cases}
\displaystyle k \sum_{i=1}^n x_i + nb = \sum_{i=1}^n y_i \\[12pt]
\displaystyle k \sum_{i=1}^n x_i^2 + b \sum_{i=1}^n x_i = \sum_{i=1}^n x_i y_i
\end{cases}
\end{align*}

Հաշվենք համակարգի որոշիչը՝
\[
\Delta =
\begin{vmatrix}
\displaystyle\sum_{i=1}^n x_i & n \\[12pt]
\displaystyle\sum_{i=1}^n x_i^2 & \displaystyle\sum_{i=1}^n x_i
\end{vmatrix}
= \left(\sum_{i=1}^n x_i\right)^2 - n \sum_{i=1}^n x_i^2
\]

$\mathbb{R}^n$ տարածությունում դիտարկենք $(1; 1; \dots; 1)^T$ և $(x_1; x_2; \dots; x_n)^T$ վեկտորները։ Կոշի–Բունյակովսկու անհավասարման համաձայն՝
\[
\left(\sum_{i=1}^n 1 \cdot x_i\right)^2 \leqslant \left(\sum_{i=1}^n 1^2\right) \left(\sum_{i=1}^n x_i^2\right) = n \sum_{i=1}^n x_i^2
\]
Հավասարման նշանը տեղի ունի այն և միայն այն դեպքում, երբ այդ վեկտորները գծորեն կախված են, այս պարագայում՝ $x_1 = x_2 = \dots = x_n$։ Քանի որ դիտարկվող խնդրում կետերի աբսցիսները բոլորը հավասար չեն (կետերը չեն գտնվում մեկ ուղղահայաց ուղղի վրա), անհավասարումը խիստ է՝
\[
\left(\sum_{i=1}^n x_i\right)^2 < n \sum_{i=1}^n x_i^2 \implies \Delta = \left(\sum_{i=1}^n x_i\right)^2 - n \sum_{i=1}^n x_i^2 < 0
\]
Հետևաբար, $\Delta \neq 0$, և համակարգն ունի միակ $(k^*, b^*)$ ստացիոնար կետը։

Այժմ օգտվելով Կրամերի կանոնից՝ գտնենք լուծումը։
\[
k = \frac{
\begin{vmatrix}
\sum_{i=1}^n y_i & n \\[6pt]
\sum_{i=1}^n x_i y_i & \sum_{i=1}^n x_i
\end{vmatrix}}{\Delta}
= \frac{\displaystyle\sum_{i=1}^n x_i \cdot \sum_{i=1}^n y_i - n \sum_{i=1}^n x_i y_i}{\displaystyle\left(\sum_{i=1}^n x_i\right)^2 - n \sum_{i=1}^n x_i^2}
\]
\[
b = \frac{1}{n} \left( \sum_{i=1}^n y_i - k \sum_{i=1}^n x_i \right) = \underbrace{\frac{\sum_{i=1}^n y_i}{n}}_{\bar{y}} - k \cdot \underbrace{\frac{\sum_{i=1}^n x_i}{n}}_{\bar{x}} = \bar{y} - k\bar{x}
\]
Այստեղից՝ $\bar{y} = k\bar{x} + b$, ինչը նշանակում է, որ $M(\bar{x}, \bar{y})$ կետը պատկանում է գծային ֆունկցիայի գրաֆիկին։

\par\medskip\noindent\textbf{Օրինակ։} Տրված են $M_1(2, 7)$, $M_2(4, 3)$, $M_3(5, 6)$, $M_4(6, 8)$ կետերը։ Կատարել գծային ապրոքսիմացիա։

\par\medskip\noindent\textbf{Լուծում։}
\begin{align*}
k &= \frac{(2+4+5+6)(7+3+6+8) - 4(14+12+30+48)}{(2+4+5+6)^2 - 4(4+16+25+36)} = \\[8pt]
&= \frac{17 \cdot 24 - 4 \cdot 104}{289 - 4 \cdot 81} = \frac{408 - 416}{289 - 324} = \frac{-8}{-35} = \frac{8}{35}
\end{align*}
\[
b = \frac{1}{4}\left(24 - \frac{8}{35} \cdot 17\right) = 6 - \frac{34}{35} = \frac{210 - 34}{35} = \frac{176}{35}
\]
Այսպիսով՝ որոնելի գծային ֆունկցիան ունի հետևյալ տեսքը՝
\[
y = \frac{8}{35}x + \frac{176}{35}
\]

\begin{figure}[htbp]
    \centering
    \begin{tikzpicture}[x=1.0cm, y=0.6cm]
        % Թեթև ցանց
        \draw[very thin, gray!25] (0,0) grid (7,8.5);

        % Առանցքներ և բաժանումներ
        \draw[->] (-0.4,0) -- (7.5,0) node[right] {$x$};
        \draw[->] (0,-0.4) -- (0,9.0) node[left] {$y$};
        \node[below left] at (0,0) {\footnotesize $0$};
        \foreach \x in {1,...,7}
            \draw (\x,0.12) -- (\x,-0.12) node[below] {\footnotesize $\x$};
        \foreach \y in {1,...,8}
            \draw (0.08,\y) -- (-0.08,\y) node[left] {\footnotesize $\y$};

        % Մոտարկող ուղիղ՝ y = (8/35)x + 176/35
        \draw[thick, blue] (0.5,5.143) -- (7.0,6.629) node[right] {$y = \frac{8}{35}x + \frac{176}{35}$};

        % Շեղումների հատվածներ (residuals)
        \draw[dashed, gray] (2,7) -- (2,5.486);
        \draw[dashed, gray] (4,3) -- (4,5.943);
        \draw[dashed, gray] (5,6) -- (5,6.171);
        \draw[dashed, gray] (6,8) -- (6,6.400);

        % Միջին կետը՝ M(4.25, 6), որն անցնում է ուղղի վրայով
        \filldraw[fill=white, draw=blue, thick] (4.25,6) circle (2.2pt) node[above left] {\small $M(\bar{x}, \bar{y})$};

        % Տրված կետերը
        \fill[red] (2,7) circle (2.2pt) node[above] {$M_1(2, 7)$};
        \fill[red] (4,3) circle (2.2pt) node[below] {$M_2(4, 3)$};
        \fill[red] (5,6) circle (2.2pt) node[below right] {$M_3(5, 6)$};
        \fill[red] (6,8) circle (2.2pt) node[above] {$M_4(6, 8)$};
    \end{tikzpicture}
    \caption{Գծային ապրոքսիմացիայի օրինակի գրաֆիկը}
\end{figure}

\section{Բազմանդամային ապրոքսիմացիա}

Դիցուք տրված են $M_1(x_1, y_1)$, $M_2(x_2, y_2)$, $\dots$, $M_n(x_n, y_n)$ կետերը։ Որպես ապրոքսիմացիա իրականացնող կոր ընտրենք $m$-րդ աստիճանի $P(x) = a_0 + a_1 x + a_2 x^2 + \dots + a_m x^m$ բազմանդամը, այնպես որ $m < n - 1$։

\begin{figure}[htbp]
    \centering
    \begin{tikzpicture}[scale=1]
        % Կոորդինատային առանցքներ
        \draw[->] (-2.4,0) -- (5.2,0) node[right] {$x$};
        \draw[->] (0,-0.5) -- (0,3.2) node[left] {$y$};

        % P(x) մոտարկող կորը (անցնում է կետերի մեջտեղով)
        \draw[thick] plot[smooth] coordinates {(-2.1,0.85) (-1.4,1.08) (0.45,1.75) (1.4,2.0) (2.4,1.8) (3.5,1.3) (4.6,1.05)};
        \node[right] at (4.6,1.05) {$P(x)$};

        % M_1 (կորից վերև)
        \fill (-1.4,1.60) circle (2pt) node[above] {$M_1$};
        \draw (-1.4,1.60) -- (-1.4,1.08);

        % M_2 (կորից ներքև)
        \fill (0.45,1.15) circle (2pt) node[right] {$M_2$};
        \draw (0.45,1.15) -- (0.45,1.75);

        % M_3 (կորից վերև)
        \fill (1.4,2.65) circle (2pt) node[above] {$M_3$};
        \draw (1.4,2.65) -- (1.4,2.0);

        % M_i և իր պրոյեկցիաները
        \fill (2.4,2.35) circle (2pt) node[above right] {$M_i$};
        \draw (2.4,2.35) -- (2.4,1.8);
        \draw[dashed] (2.4,1.8) -- (2.4,0) node[below] {$x_i$};
        \draw[dashed] (2.4,2.35) -- (0,2.35) node[left] {$y_i$};

        % M_n (կորից ներքև)
        \fill (3.5,0.75) circle (2pt) node[below] {$M_n$};
        \draw (3.5,0.75) -- (3.5,1.3);
    \end{tikzpicture}
    \caption{Բազմանդամային ապրոքսիմացիա}
\end{figure}

Կիրառենք փոքրագույն քառակուսիների մեթոդը՝
\[
S = S(a_0, a_1, \dots, a_m) = \sum_{i=1}^n \bigl(y_i - (a_0 + a_1 x_i + a_2 x_i^2 + \dots + a_m x_i^m)\bigr)^2 \to \min
\]

Բազմանդամի $m$ աստիճանը մեծացնելիս $S$-ը չի աճում, և սահմանային $m = n - 1$ դեպքում բազմանդամը վերածվում է ինտերպոլյացիոն կորի՝ $S = 0$ արժեքով։ Սակայն եթե տվյալները պարունակում են աղմուկ, չափազանց բարձր $m$-ի ընտրությունը հանգեցնում է \textit{գերհարմարեցման}. կորը ճշգրիտ անցնում է աղմկոտ կետերով, սակայն կետերի միջև առաջացնում է կտրուկ տատանումներ և կորցնում կանխատեսողական հատկությունը։ Այդ պատճառով գործնականում ընտրում են $m < n - 1$։

Հաշվենք $S$ ֆունկցիայի մասնակի ածանցյալները՝
\begin{align*}
\frac{\partial S}{\partial a_0} &= -2 \sum_{i=1}^n \bigl(y_i - (a_0 + a_1 x_i + a_2 x_i^2 + \dots + a_m x_i^m)\bigr) \\
\frac{\partial S}{\partial a_1} &= -2 \sum_{i=1}^n \bigl(y_i - (a_0 + a_1 x_i + a_2 x_i^2 + \dots + a_m x_i^m)\bigr) x_i \\
\frac{\partial S}{\partial a_2} &= -2 \sum_{i=1}^n \bigl(y_i - (a_0 + a_1 x_i + a_2 x_i^2 + \dots + a_m x_i^m)\bigr) x_i^2 \\
&\;\;\vdots \\
\frac{\partial S}{\partial a_m} &= -2 \sum_{i=1}^n \bigl(y_i - (a_0 + a_1 x_i + a_2 x_i^2 + \dots + a_m x_i^m)\bigr) x_i^m
\end{align*}

Գտնենք $S$ ֆունկցիայի ստացիոնար կետը, որտեղ և ֆունկցիան կընդունի իր փոքրագույն արժեքը՝

\begin{align}
\begin{cases}
\dfrac{\partial S}{\partial a_0} = 0 \\[8pt]
\dfrac{\partial S}{\partial a_1} = 0 \\[4pt]
\;\;\vdots \\[4pt]
\dfrac{\partial S}{\partial a_m} = 0
\end{cases}
&\Longleftrightarrow
\begin{cases}
\displaystyle a_0 n + a_1 \sum_{i=1}^n x_i + a_2 \sum_{i=1}^n x_i^2 + \dots + a_m \sum_{i=1}^n x_i^m = \sum_{i=1}^n y_i \\[14pt]
\displaystyle a_0 \sum_{i=1}^n x_i + a_1 \sum_{i=1}^n x_i^2 + a_2 \sum_{i=1}^n x_i^3 + \dots + a_m \sum_{i=1}^n x_i^{m+1} = \sum_{i=1}^n x_i y_i \\[8pt]
\;\;\vdots \\[8pt]
\displaystyle a_0 \sum_{i=1}^n x_i^m + a_1 \sum_{i=1}^n x_i^{m+1} + \dots + a_m \sum_{i=1}^n x_i^{2m} = \sum_{i=1}^n x_i^m y_i
\end{cases}
\label{eq:poly_system}
\end{align}

Դիտարկենք հետևյալ մատրիցը՝
\[
A =
\begin{pmatrix}
1 & x_1 & x_1^2 & \dots & x_1^m \\
1 & x_2 & x_2^2 & \dots & x_2^m \\
1 & x_3 & x_3^2 & \dots & x_3^m \\
\vdots & \vdots & \vdots & \ddots & \vdots \\
1 & x_n & x_n^2 & \dots & x_n^m
\end{pmatrix}_{n \times (m+1)} \quad \text{(Վանդերմոնդի մատրից)}
\]
և սյուն-վեկտորները՝
\[
X =
\begin{pmatrix}
a_0 \\
a_1 \\
\vdots \\
a_m
\end{pmatrix}_{(m+1) \times 1},
\qquad
Y =
\begin{pmatrix}
y_1 \\
y_2 \\
\vdots \\
y_n
\end{pmatrix}_{n \times 1}
\]

Դիտարկենք հետևյալ արտադրյալները՝
\begin{align*}
A^T \cdot A &=
\begin{pmatrix}
1 & 1 & 1 & \dots & 1 \\
x_1 & x_2 & x_3 & \dots & x_n \\
x_1^2 & x_2^2 & x_3^2 & \dots & x_n^2 \\
\vdots & \vdots & \vdots & \ddots & \vdots \\
x_1^m & x_2^m & x_3^m & \dots & x_n^m
\end{pmatrix}
\cdot
\begin{pmatrix}
1 & x_1 & x_1^2 & \dots & x_1^m \\
1 & x_2 & x_2^2 & \dots & x_2^m \\
1 & x_3 & x_3^2 & \dots & x_3^m \\
\vdots & \vdots & \vdots & \ddots & \vdots \\
1 & x_n & x_n^2 & \dots & x_n^m
\end{pmatrix}_{n \times (m+1)} \\[12pt]
&=
\begin{pmatrix}
n & \displaystyle\sum_{i=1}^n x_i & \displaystyle\sum_{i=1}^n x_i^2 & \dots & \displaystyle\sum_{i=1}^n x_i^m \\[14pt]
\displaystyle\sum_{i=1}^n x_i & \displaystyle\sum_{i=1}^n x_i^2 & \displaystyle\sum_{i=1}^n x_i^3 & \dots & \displaystyle\sum_{i=1}^n x_i^{m+1} \\[12pt]
\vdots & \vdots & \vdots & \ddots & \vdots \\[8pt]
\displaystyle\sum_{i=1}^n x_i^m & \displaystyle\sum_{i=1}^n x_i^{m+1} & \displaystyle\sum_{i=1}^n x_i^{m+2} & \dots & \displaystyle\sum_{i=1}^n x_i^{2m}
\end{pmatrix}_{(m+1) \times (m+1)}
\end{align*}

\begin{align}
(A^T \cdot A) \cdot X &=
\begin{pmatrix}
n & \sum_{i=1}^n x_i & \sum_{i=1}^n x_i^2 & \dots & \sum_{i=1}^n x_i^m \\[6pt]
\sum_{i=1}^n x_i & \sum_{i=1}^n x_i^2 & \sum_{i=1}^n x_i^3 & \dots & \sum_{i=1}^n x_i^{m+1} \\[6pt]
\sum_{i=1}^n x_i^2 & \sum_{i=1}^n x_i^3 & \sum_{i=1}^n x_i^4 & \dots & \sum_{i=1}^n x_i^{m+2} \\[4pt]
\vdots & \vdots & \vdots & \ddots & \vdots \\[4pt]
\sum_{i=1}^n x_i^m & \sum_{i=1}^n x_i^{m+1} & \sum_{i=1}^n x_i^{m+2} & \dots & \sum_{i=1}^n x_i^{2m}
\end{pmatrix}_{(m+1) \times (m+1)}
\cdot
\begin{pmatrix}
a_0 \\
a_1 \\
a_2 \\
\vdots \\
a_m
\end{pmatrix}_{(m+1) \times 1} \nonumber \\[12pt]
\end{align}

\begin{align}
&=
\begin{pmatrix}
a_0 n + a_1 \sum_{i=1}^n x_i + a_2 \sum_{i=1}^n x_i^2 + \dots + a_m \sum_{i=1}^n x_i^m \\[6pt]
a_0 \sum_{i=1}^n x_i + a_1 \sum_{i=1}^n x_i^2 + a_2 \sum_{i=1}^n x_i^3 + \dots + a_m \sum_{i=1}^n x_i^{m+1} \\[6pt]
a_0 \sum_{i=1}^n x_i^2 + a_1 \sum_{i=1}^n x_i^3 + a_2 \sum_{i=1}^n x_i^4 + \dots + a_m \sum_{i=1}^n x_i^{m+2} \\[4pt]
\vdots \\[4pt]
a_0 \sum_{i=1}^n x_i^m + a_1 \sum_{i=1}^n x_i^{m+1} + a_2 \sum_{i=1}^n x_i^{m+2} + \dots + a_m \sum_{i=1}^n x_i^{2m}
\end{pmatrix}_{(m+1) \times 1}
\label{eq:ata_x}
\end{align}

\begin{align}
A^T \cdot Y &=
\begin{pmatrix}
1 & 1 & 1 & \dots & 1 \\
x_1 & x_2 & x_3 & \dots & x_n \\
x_1^2 & x_2^2 & x_3^2 & \dots & x_n^2 \\
\vdots & \vdots & \vdots & \ddots & \vdots \\
x_1^m & x_2^m & x_3^m & \dots & x_n^m
\end{pmatrix}_{(m+1) \times n}
\cdot
\begin{pmatrix}
y_1 \\
y_2 \\
y_3 \\
\vdots \\
y_n
\end{pmatrix}_{n \times 1} \nonumber \\[12pt]
&=
\begin{pmatrix}
\sum_{i=1}^n y_i \\[6pt]
\sum_{i=1}^n x_i y_i \\[6pt]
\sum_{i=1}^n x_i^2 y_i \\[4pt]
\vdots \\[4pt]
\sum_{i=1}^n x_i^m y_i
\end{pmatrix}_{(m+1) \times 1}
\label{eq:aty}
\end{align}

Հաշվի առնելով \eqref{eq:ata_x}-ը և \eqref{eq:aty}-ը՝ \eqref{eq:poly_system} համակարգը կարելի է գրել հետևյալ մատրիցային հավասարման միջոցով՝
\[
(A^T \cdot A) \cdot X = A^T \cdot Y
\]
\[
\boxed{X = (A^T \cdot A)^{-1} \cdot A^T \cdot Y}
\]

\par\medskip\noindent\textbf{Դիտողություն (լուծման միակության մասին)։} Եթե կետերի $x_1, \dots, x_n$ աբսցիսները զույգ առ զույգ տարբեր են և $m < n$, ապա առավելագույնը $m$-րդ աստիճանի ոչ զրոյական $P(x)$ բազմանդամը չի կարող զրոյանալ $n > m$ տարբեր կետերում, որտեղից հետևում է, որ $A$ մատրիցի սյուները գծորեն անկախ են, $\det(A^T A) \neq 0$, և նորմալ համակարգն ունի միակ $X^*$ լուծումը։ Քանի որ $S = \sum_{i=1}^n \bigl(y_i - P(x_i)\bigr)^2 \geqslant 0$-ն անընդհատ և դիֆերենցելի ֆունկցիա է, և $\|X\| \to +\infty$ դեպքում $S \to +\infty$, ապա 1-ին գլխի համանմանությամբ գտնված միակ $X^*$ ստացիոնար կետում $S$-ն ընդունում է իր փոքրագույն արժեքը։
\par\medskip\noindent\textbf{Օրինակ։} Տրված են $M_1(2, 7)$, $M_2(4, 3)$, $M_3(5, 6)$, $M_4(6, 8)$ կետերը։

\par\medskip\noindent\textbf{Լուծում։} Այստեղ $n = 4$։ Պայմանից ունենք՝
\[
m < n - 1 \implies m < 3 \implies m = 2 \implies P(x) = a_0 + a_1 x + a_2 x^2
\]
Կազմենք $A$, $A^T$ մատրիցները և $Y$, $X$ սյուն-վեկտորները՝
\[
A =
\begin{pmatrix}
1 & x_1 & x_1^2 \\
1 & x_2 & x_2^2 \\
1 & x_3 & x_3^2 \\
1 & x_4 & x_4^2
\end{pmatrix}_{4 \times 3}
=
\begin{pmatrix}
1 & 2 & 4 \\
1 & 4 & 16 \\
1 & 5 & 25 \\
1 & 6 & 36
\end{pmatrix}
\]
\[
A^T =
\begin{pmatrix}
1 & 1 & 1 & 1 \\
2 & 4 & 5 & 6 \\
4 & 16 & 25 & 36
\end{pmatrix}_{3 \times 4}
\]
\[
Y =
\begin{pmatrix}
7 \\
3 \\
6 \\
8
\end{pmatrix},
\qquad
X =
\begin{pmatrix}
a_0 \\
a_1 \\
a_2
\end{pmatrix}
\]
Քանի որ $X = (A^T \cdot A)^{-1} \cdot A^T \cdot Y$, հաշվենք $A^T \cdot A$ մատրիցը՝
\[
A^T \cdot A =
\begin{pmatrix}
4 & 17 & 81 \\
17 & 81 & 413 \\
81 & 413 & 2193
\end{pmatrix}
\]

\[
(A^T A)^{-1} = \frac{C^T}{\Delta(A^T A)}
\]
\begin{align*}
\Delta(A^T A) &=
\begin{vmatrix}
4 & 17 & 81 \\
17 & 81 & 413 \\
81 & 413 & 2193
\end{vmatrix} \\[8pt]
&= 4
\begin{vmatrix}
81 & 413 \\
413 & 2193
\end{vmatrix}
- 17
\begin{vmatrix}
17 & 413 \\
81 & 2193
\end{vmatrix}
+ 81
\begin{vmatrix}
17 & 81 \\
81 & 413
\end{vmatrix} \\[8pt]
&= 28256 - 65076 + 37260 = 440
\end{align*}

Հաշվենք $A^T A$ մատրիցի տարրերի հանրահաշվական լրացումները՝
\begin{align*}
C_{11} &=
\begin{vmatrix}
81 & 413 \\
413 & 2193
\end{vmatrix}
= 7064 \\[6pt]
C_{12} &= -
\begin{vmatrix}
17 & 413 \\
81 & 2193
\end{vmatrix}
= -3828 \\[6pt]
C_{13} &=
\begin{vmatrix}
17 & 81 \\
81 & 413
\end{vmatrix}
= 460 \\[6pt]
C_{21} &= -
\begin{vmatrix}
17 & 81 \\
413 & 2193
\end{vmatrix}
= -3828 \\[6pt]
C_{22} &=
\begin{vmatrix}
4 & 81 \\
81 & 2193
\end{vmatrix}
= 2211 \\[6pt]
C_{23} &= -
\begin{vmatrix}
4 & 17 \\
81 & 413
\end{vmatrix}
= -275 \\[6pt]
C_{31} &=
\begin{vmatrix}
17 & 81 \\
81 & 413
\end{vmatrix}
= 460 \\[6pt]
C_{32} &= -
\begin{vmatrix}
4 & 81 \\
17 & 413
\end{vmatrix}
= -275 \\[6pt]
C_{33} &=
\begin{vmatrix}
4 & 17 \\
17 & 81
\end{vmatrix}
= 35
\end{align*}

Հետևաբար՝
\[
C =
\begin{pmatrix}
7064 & -3828 & 460 \\
-3828 & 2211 & -275 \\
460 & -275 & 35
\end{pmatrix}
\]
\[
(A^T A)^{-1} = \frac{1}{440}
\begin{pmatrix}
7064 & -3828 & 460 \\
-3828 & 2211 & -275 \\
460 & -275 & 35
\end{pmatrix}
=
\begin{pmatrix}
\dfrac{883}{55} & -\dfrac{87}{10} & \dfrac{23}{22} \\[10pt]
-\dfrac{87}{10} & \dfrac{201}{40} & -\dfrac{5}{8} \\[10pt]
\dfrac{23}{22} & -\dfrac{5}{8} & \dfrac{7}{88}
\end{pmatrix}
\]

Գտնենք $X$ վեկտորը՝
\begin{align*}
X &=
\begin{pmatrix}
\dfrac{883}{55} & -\dfrac{87}{10} & \dfrac{23}{22} \\[10pt]
-\dfrac{87}{10} & \dfrac{201}{40} & -\dfrac{5}{8} \\[10pt]
\dfrac{23}{22} & -\dfrac{5}{8} & \dfrac{7}{88}
\end{pmatrix}
\cdot
\begin{pmatrix}
1 & 1 & 1 & 1 \\
2 & 4 & 5 & 6 \\
4 & 16 & 25 & 36
\end{pmatrix}
\cdot
\begin{pmatrix}
7 \\
3 \\
6 \\
8
\end{pmatrix} \\[12pt]
&=
\begin{pmatrix}
\dfrac{883}{55} & -\dfrac{87}{10} & \dfrac{23}{22} \\[10pt]
-\dfrac{87}{10} & \dfrac{201}{40} & -\dfrac{5}{8} \\[10pt]
\dfrac{23}{22} & -\dfrac{5}{8} & \dfrac{7}{88}
\end{pmatrix}
\cdot
\begin{pmatrix}
24 \\
104 \\
514
\end{pmatrix} \\[12pt]
&=
\begin{pmatrix}
\dfrac{21192}{55} - \dfrac{9048}{10} + \dfrac{11822}{22} \\[10pt]
-\dfrac{1044}{5} + \dfrac{2613}{5} - \dfrac{1285}{4} \\[10pt]
\dfrac{276}{11} - 65 + \dfrac{1799}{44}
\end{pmatrix}
=
\begin{pmatrix}
\dfrac{983}{55} \\[10pt]
-\dfrac{149}{20} \\[10pt]
\dfrac{43}{44}
\end{pmatrix}
\end{align*}

Այսպիսով՝ որոնելի բազմանդամն ունի հետևյալ տեսքը՝
\[
\boxed{P(x) = \frac{983}{55} - \frac{149}{20}x + \frac{43}{44}x^2}
\]

\begin{figure}[htbp]
    \centering
    \begin{tikzpicture}[x=1.0cm, y=0.58cm]
        % Թեթև ցանց
        \draw[very thin, gray!25] (0,0) grid (7,9.5);

        % Առանցքներ և բաժանումներ
        \draw[->] (-0.4,0) -- (7.5,0) node[right] {$x$};
        \draw[->] (0,-0.4) -- (0,10.0) node[left] {$y$};
        \node[below left] at (0,0) {\footnotesize $0$};
        \foreach \x in {1,...,7}
            \draw (\x,0.12) -- (\x,-0.12) node[below] {\footnotesize $\x$};
        \foreach \y in {1,...,9}
            \draw (0.08,\y) -- (-0.08,\y) node[left] {\footnotesize $\y$};

        % Մոտարկող պարաբոլը՝ P(x) = 983/55 - (149/20)x + (43/44)x^2
        \draw[thick, blue] plot[smooth] coordinates {
            (1.45,9.13) (1.75,7.83) (2.0,6.88) (2.5,5.36) (3.0,4.32)
            (3.5,3.77) (3.81,3.67) (4.0,3.71) (4.5,4.14) (5.0,5.05)
            (5.5,6.46) (6.0,8.35) (6.25,9.49)
        };
        \node[right, blue] at (6.25,9.35) {$P(x) = \frac{983}{55} - \frac{149}{20}x + \frac{43}{44}x^2$};

        % Շեղումների հատվածներ (residuals)
        \draw[dashed, gray] (2,7) -- (2,6.88);
        \draw[dashed, gray] (4,3) -- (4,3.71);
        \draw[dashed, gray] (5,6) -- (5,5.05);
        \draw[dashed, gray] (6,8) -- (6,8.35);

        % Տրված կետերը
        \fill[red] (2,7) circle (2.2pt) node[right] {$M_1(2, 7)$};
        \fill[red] (4,3) circle (2.2pt) node[below] {$M_2(4, 3)$};
        \fill[red] (5,6) circle (2.2pt) node[above left] {$M_3(5, 6)$};
        \fill[red] (6,8) circle (2.2pt) node[right] {$M_4(6, 8)$};
    \end{tikzpicture}
    \caption{Քառակուսային ապրոքսիմացիայի օրինակի գրաֆիկը}
\end{figure}

\subsection{Որոշ ոչ գծային ֆունկցիաների գծայնացումը}

Գործնականում հաճախ հանդիպում են ոչ գծային կախվածություններ, որոնք պարզ փոփոխականների փոխարինման միջոցով կարելի է բերել գծային փոքրագույն քառակուսիների խնդրի.
\begin{itemize}
    \item \textbf{Ցուցչային ֆունկցիա՝} $y = a e^{bx}$ ($a > 0$)։ Լոգարիթմելով երկու կողմը՝ ստանում ենք $\ln y = \ln a + bx$։ Նշանակելով $\tilde{y}_i = \ln y_i$ և $\tilde{b} = \ln a$՝ ունենում ենք գծային ֆունկցիա $\tilde{y} = bx + \tilde{b}$, որից գտնում ենք $b$-ն և $a = e^{\tilde{b}}$-ն։
    \item \textbf{Աստիճանային ֆունկցիա՝} $y = a x^b$ ($a > 0, \, x > 0$)։ Լոգարիթմելով՝ $\ln y = \ln a + b \ln x$։ Կատարելով $\tilde{y}_i = \ln y_i$, $\tilde{x}_i = \ln x_i$ և $\tilde{b} = \ln a$ նշանակումները՝ խնդիրը կրկին բերվում է $\tilde{y} = b \tilde{x} + \tilde{b}$ գծային ապրոքսիմացիայի։
    \item \textbf{Լոգարիթմական ֆունկցիա՝} $y = a \ln x + b$ ($x > 0$)։ Նշանակելով $\tilde{x}_i = \ln x_i$՝ անմիջապես ստանում ենք գծային ֆունկցիա՝ $y = a \tilde{x} + b$։
\end{itemize}

\section{Մոտարկման որակի գնահատումը}

Տարբեր աստիճանների մոտարկող բազմանդամների ճշգրտությունը համեմատելու համար կիրառվում են մի քանի հիմնական թվային չափանիշներ.
\begin{enumerate}
    \item \textbf{Միջին քառակուսային սխալ (MSE) և միջին քառակուսային շեղում (RMSE).}
    \[
    \mathrm{MSE} = \frac{1}{n} \sum_{i=1}^n \bigl(y_i - P(x_i)\bigr)^2 = \frac{S}{n}, \qquad \mathrm{RMSE} = \sqrt{\mathrm{MSE}} = \sqrt{\frac{S}{n}}
    \]
    Որքան փոքր է $\mathrm{RMSE}$-ն, այնքան ավելի մոտ է կառուցված կորը տրված կետերին։
    \item \textbf{Դետերմինացիայի գործակից ($R^2$)։} Այս ցուցանիշը գնահատում է ապրոքսիմացիայի ճշգրտությունը և ցույց է տալիս, թե մոտարկող ֆունկցիան որքանով է նվազեցնում շեղումը՝ համեմատած $y_i$ արժեքների $\bar{y}$ միջինից ունեցած ընդհանուր շեղման հետ՝
    \[
    R^2 = 1 - \frac{\sum_{i=1}^n \bigl(y_i - P(x_i)\bigr)^2}{\sum_{i=1}^n (y_i - \bar{y})^2} = 1 - \frac{S}{\mathrm{SST}}
    \]
    որտեղ $\bar{y} = \dfrac{1}{n}\sum_{i=1}^n y_i$-ն միջին արժեքն է, իսկ $\mathrm{SST} = \sum_{i=1}^n (y_i - \bar{y})^2$-ը՝ ընդհանուր շեղումների քառակուսիների գումարը։ Սովորաբար $0 \leqslant R^2 \leqslant 1$, և որքան $R^2$-ը մոտ է $1$-ին, այնքան բարձր է մոտարկման որակը։
\end{enumerate}

Համեմատենք նախորդ գլուխներում նույն $M_1(2, 7)$, $M_2(4, 3)$, $M_3(5, 6)$, $M_4(6, 8)$ կետերի համար ստացված գծային ($m = 1$) և քառակուսային ($m = 2$) մոտարկող ֆունկցիաները։ Այստեղ $\bar{y} = 6$ և $\mathrm{SST} = (7-6)^2 + (3-6)^2 + (6-6)^2 + (8-6)^2 = 14$։ Հաշվարկների արդյունքներն ամփոփված են Աղյուսակ 1-ում։

\begin{table}[htbp]
    \centering
    \small
    \renewcommand{\arraystretch}{1.45}
    \caption{Գծային և քառակուսային մոտարկումների սխալների համեմատություն}
    \begin{tabular}{@{}p{4.3cm}cccc@{}}
        \toprule
        \textbf{Մոտարկող ֆունկցիան} & $S$ & $\mathrm{MSE}$ & $\mathrm{RMSE}$ & $R^2$ \tabularnewline
        \midrule
        \raggedright Գծային ($m = 1$) & $13.5429$ & $3.3857$ & $1.8400$ & $0.0327$ \tabularnewline
        \raggedright Քառակուսային ($m = 2$) & $1.5364$ & $0.3841$ & $0.6198$ & $0.8903$ \tabularnewline
        \bottomrule
    \end{tabular}
\end{table}

Ինչպես երևում է աղյուսակից, գծային մոտարկումը վատ է արտացոլում տրված կետերի ոչ գծային բնույթը ($R^2 \approx 3.3\%$), մինչդեռ $m = 2$ բազմանդամը նվազեցնում է քառակուսային սխալը գրեթե $9$ անգամ՝ ապահովելով $R^2 \approx 89.0\%$։
\newpage
\section*{Եզրակացություն}
\phantomsection
\addcontentsline{toc}{section}{\texorpdfstring{\MakeUppercase{Եզրակացություն}}{Եզրակացություն}}

Սույն աշխատանքում ուսումնասիրվել են փոքրագույն քառակուսիների մեթոդի տեսական հիմունքները և դրա կիրառությունը գծային ու բազմանդամային ապրոքսիմացիայի խնդիրներում։ Աշխատանքի ընթացքում կատարվել են հետևյալ հիմնական քայլերն ու ստացվել են հետևյալ արդյունքները.

\begin{enumerate}
    \item Սահմանվել է փորձարարական կետերի հավաքածուի մոտարկման խնդիրը և ցույց է տրվել ապրոքսիմացիայի հիմնարար տարբերությունը ինտերպոլյացիայից՝ աղմկոտ տվյալների պայմաններում։
    \item Գծային ապրոքսիմացիայի ($y = kx + b$) համար արտածվել է նորմալ հավասարումների համակարգը, Կոշի–Բունյակովսկու անհավասարմամբ ապացուցվել է $\Delta < 0$ լուծելիության պայմանը, և Կրամերի կանոնով ստացվել են $k$ ու $b$ գործակիցների բացահայտ բանաձևերը։
    \item Բազմանդամային ապրոքսիմացիայի ($P(x) = a_0 + a_1 x + \dots + a_m x^m$) համար կառուցվել է նորմալ հավասարումների համակարգը և ներկայացվել է մատրիցային տեսքով՝ $(A^T A)X = A^T Y \implies X = (A^T A)^{-1} A^T Y$։
    \item Կոնկրետ թվային օրինակների վրա կառուցվել են գծային ($m = 1$) և քառակուսային ($m = 2$) մոտարկող ֆունկցիաները, և $\mathrm{MSE}$, $\mathrm{RMSE}$ ու $R^2$ չափանիշների հիման վրա համեմատվել է դրանց ճշգրտությունը (քառակուսային մոտարկման դեպքում դետերմինացիայի գործակիցը $R^2 = 0.0327$-ից աճել է մինչև $R^2 = 0.8903$)։
    \item Դիտարկվել են նաև բազմանդամի աստիճանի ընտրությանը վերաբերող գերհարմարեցման (overfitting) ռիսկերը և ցուցադրվել է, թե ինչպես կարելի է մի շարք ոչ գծային կախվածություններ (ցուցչային, աստիճանային, լոգարիթմական) բերել գծային փոքրագույն քառակուսիների խնդրի։
\end{enumerate}

Ամփոփելով՝ կարելի է պնդել, որ փոքրագույն քառակուսիների մեթոդը արդյունավետ և ունիվերսալ գործիք է փորձարարական տվյալների մոտարկման և կախվածությունների վերլուծության համար, իսկ մատրիցային ձևակերպումը հիմք է ստեղծում ավելի բարդ թվային մեթոդների և մատրիցային վերլուծությունների կիրառման համար։

\newpage 

\renewcommand{\refname}{ԳՐԱԿԱՆՈՒԹՅԱՆ ՑԱՆԿ}
\begin{thebibliography}{9}

\bibitem{hakobyan1}
Հակոբյան Յու. Ռ. \textit{Թվային մեթոդներ, մաս I}. — Երևան. Հանրագիտարան «Արմենիկա», 2003. — 222 էջ։

\bibitem{hakobyan2}
Հակոբյան Յու. Ռ. \textit{Թվային մեթոդներ, մաս II}. — Երևան. «ՎՄՎ-Պրինտ», 2007. — 246 էջ։

\bibitem{bakhvalov}
Бахвалов Н. С., Жидков Н. П., Кобельков Г. М. \textit{Численные методы}. — 9-е изд. — М.. «Лаборатория знаний», 2020. — 636 с.

\bibitem{samarskii}
Самарский А. А., Гулин А. В. \textit{Численные методы}. — М. «Наука», 1989. — 432 с.

\bibitem{burden}
Burden R. L., Faires J. D. \textit{Numerical Analysis}. — 9th ed. — Boston, MA. Brooks/Cole, Cengage Learning, 2011. — 872 p.

\bibitem{strang}
Strang G. \textit{Introduction to Linear Algebra}. — 5th ed. — Wellesley, MA. Wellesley-Cambridge Press, 2016. — 584 p.

\end{thebibliography}

\end{document}